% ===== PRÉAMBULE CANONIQUE — à coller VERBATIM en tête de CHAQUE .tex =====
\documentclass[11pt,a4paper]{article}
\usepackage[utf8]{inputenc}\usepackage[T1]{fontenc}\usepackage{lmodern}
\usepackage[french]{babel}
\usepackage{amsmath,amssymb,mathtools}\usepackage{icomma}
\usepackage{siunitx}\sisetup{locale=FR}  % \qty{}{}, \si{}, \ang{}
\usepackage{xcolor}\usepackage{colortbl}
\usepackage{tikz}\usepackage{pgfplots}\pgfplotsset{compat=1.18}
\usepackage[siunitx,europeanresistors]{circuitikz} % circuits (élec.)
\usepackage[most]{tcolorbox}\usepackage{enumitem}\usepackage{array,booktabs}
\usepackage[a4paper,margin=2.2cm]{geometry}
\usetikzlibrary{arrows.meta,calc,angles,quotes,positioning,patterns,%
  backgrounds,shapes.geometric,decorations.pathmorphing,decorations.markings,intersections}
\definecolor{primary}{RGB}{222,49,99}\definecolor{primarydark}{RGB}{150,22,55}
\definecolor{defcol}{RGB}{0,114,178}\definecolor{methcol}{RGB}{52,92,148}
\definecolor{excol}{RGB}{0,135,90}\definecolor{attncol}{RGB}{200,40,55}
\tcbset{boxbase/.style={enhanced,breakable,boxrule=0.9pt,arc=2.5pt,
  left=3mm,right=3mm,top=2.3mm,bottom=2.3mm,fonttitle=\bfseries,coltitle=white}}
% --- boîtes TUTORIEL (noms EXACTS, ne pas renommer) ---
\newtcolorbox{cle}[1][]{boxbase,colback=defcol!6,colframe=defcol,
  title={\textsc{À retenir}\ifx\relax#1\relax\relax\else\textmd{ : \itshape #1}\fi}}
\newtcolorbox{methode}[1][]{boxbase,colback=methcol!7,colframe=methcol,sharp corners,
  title={\textsc{Méthode}\ifx\relax#1\relax\relax\else\textmd{ --- \itshape #1}\fi}}
\newtcolorbox{exemplebox}[1][]{boxbase,colback=excol!5,colframe=excol,
  title={\textsc{Exemple}\ifx\relax#1\relax\relax\else\textmd{ : \itshape #1}\fi}}
\newtcolorbox{piege}[1][]{boxbase,colback=attncol!6,colframe=attncol,
  title={\textsc{Piège}\ifx\relax#1\relax\relax\else\textmd{ : \itshape #1}\fi}}
\newtcolorbox{remarque}[1][]{boxbase,colback=black!4,colframe=black!45,
  title={\textsc{Remarque}\ifx\relax#1\relax\relax\else\textmd{ : \itshape #1}\fi}}
% --- boîtes EXERCICE / CORRIGÉ (noms EXACTS, auto-comptées) ---
\newtcolorbox[auto counter]{exo}[1][]{boxbase,colback=primary!4,colframe=primary,
  title={\textsc{Exercice}~\thetcbcounter\ifx\relax#1\relax\relax\else\textmd{ --- \itshape #1}\fi}}
\newtcolorbox[auto counter]{corrige}[1][]{boxbase,colback=excol!4,colframe=excol,
  title={\textsc{Corrigé}~\thetcbcounter\ifx\relax#1\relax\relax\else\textmd{ --- \itshape #1}\fi}}
\newcommand{\vv}[1]{\vec{#1}}\newcommand{\ds}{\displaystyle}
\newcommand{\R}{\mathbb{R}}\newcommand{\N}{\mathbb{N}}\newcommand{\Z}{\mathbb{Z}}
% =========================================================================
\begin{document}
\usetikzlibrary{babel}
\begin{exo}[résistance équivalente en série]
Trois résistances $R_1=\qty{10}{\ohm}$, $R_2=\qty{20}{\ohm}$ et $R_3=\qty{30}{\ohm}$ sont en série.
Calcule la résistance équivalente $R_{\text{eq}}$.
\begin{center}
\begin{circuitikz}[scale=0.8,transform shape]
  \draw (0,0) to[short,o-] (0.5,0) to[R,l=$R_1$] (2.2,0) to[R,l=$R_2$] (3.9,0)
        to[R,l=$R_3$] (5.6,0) to[short,-o] (6.1,0);
\end{circuitikz}
\end{center}
\end{exo}

\begin{exo}[deux résistances en parallèle]
Deux résistances $R_1=\qty{6}{\ohm}$ et $R_2=\qty{3}{\ohm}$ sont en parallèle. Calcule
$R_{\text{eq}}$ (astuce : produit sur somme).
\begin{center}
\begin{circuitikz}[scale=0.8,transform shape]
  \coordinate (A) at (0,0); \coordinate (B) at (3.6,0);
  \draw[thick] (A) -- (0,1.1); \draw (0,1.1) to[R,l=$R_1$] (3.6,1.1); \draw[thick] (3.6,1.1)--(B);
  \draw[thick] (A) -- (0,-1.1); \draw (0,-1.1) to[R,l=$R_2$] (3.6,-1.1); \draw[thick] (3.6,-1.1)--(B);
  \draw[thick] (A)--(-0.6,0); \draw[thick] (B)--(4.2,0);
\end{circuitikz}
\end{center}
\end{exo}

\begin{exo}[$n$ résistances égales en parallèle]
Trois résistances identiques $R=\qty{30}{\ohm}$ sont montées en parallèle. Que vaut $R_{\text{eq}}$ ?
(astuce : $R/n$).
\end{exo}

\begin{exo}[générateur réel]
Une pile a une f.é.m. $E=\qty{9}{\volt}$ et une résistance interne $r=\qty{1}{\ohm}$.
\begin{enumerate}[label=\alph*)]
  \item Quelle est sa tension à vide (lorsqu'elle ne débite aucun courant) ?
  \item Quelle tension fournit-elle lorsqu'elle débite $I=\qty{2}{\ampere}$ ?
\end{enumerate}
\end{exo}

\begin{exo}[intensité dans un circuit série]
Deux résistances $R_1=\qty{10}{\ohm}$ et $R_2=\qty{20}{\ohm}$ en série sont soumises à une tension
$U=\qty{30}{\volt}$. Calcule l'intensité $I$ qui les traverse.
\end{exo}

\end{document}
